\aat{This page will be used for additional experiments for the camera-ready version. Will be removed after finalizing everything. }

\begin{itemize}
    \item We need to add results for more data
    \item We need to add results for cPCA 
    \item We need to add a comparison between mean pooling and the last token 
    \item We need to add confidence level (range value)
    \item Comparison between ambiguity, Clear SP, Clear SC.
\end{itemize}

\paragraph{Validation on an Expanded Dataset.}
To test whether the conclusions established in the original local
semantic neighborhoods hold across a broader range of inputs, we
construct an additional dataset containing 1,400 semantic groups and
57,400 sentences across 13 domains. Each group contains one anchor,
20 semantic-preserving (SP) paraphrases, 10 soft semantic-changing
(SC) variants, and 10 hard SC variants. SP paraphrases vary grammatical
structure and writing style while preserving meaning. Soft SC variants
modify an attribute or semantic slot within the same topic, whereas
hard SC variants introduce a domain- or topic-level semantic change.
SC examples include edit-type and source/target-domain metadata for
per-perturbation analysis. This expanded dataset supplements the
original local analyses and evaluates whether the observed
semantic--nuisance separation extends to more diverse semantic
neighborhoods.

\paragraph{Multi-hop Reasoning Validation.}
Multi-hop reasoning requires combining multiple facts or intermediate
inferences to reach a conclusion, rather than relying on a single
explicit fact. We additionally evaluate 200 multi-hop reasoning
examples to test whether the invariant decomposition remains
selective when semantic differences require reasoning beyond direct
lexical or attribute substitutions. Table~\ref{tab:multihop_perturbation_examples} showed some example of dataset. Using the same decomposition and
intervention protocol, we compare the effects of replacing semantic
and nuisance components. As shown in
Table~\ref{tab:multihop_validation}, semantic interventions produce
larger MSE and KL changes than nuisance interventions at all reported
layers. These results support the conclusion that the discovered
invariant directions remain semantically selective in
reasoning-intensive settings.

\begin{table}[t]
\centering
\small
\caption{Illustrative multi-hop examples with nuisance (SP) and
semantic (SC) changes. SP variants preserve the reasoning and answer,
whereas SC variants change a fact or relationship and alter the answer.
The question remains fixed within each example.}
\label{tab:multihop_perturbation_examples}
\begin{tabular}{p{0.13\linewidth} p{0.65\linewidth} p{0.13\linewidth}}
\toprule
Type & Supporting Facts & Answer \\
\midrule
\multicolumn{3}{p{0.95\linewidth}}{
\textbf{Example 1:} In which country does Maya live?} \\
Anchor &
Maya lives in Brookfield. Brookfield is in Canada. &
Canada \\
SP &
Maya resides in Brookfield. Brookfield is located in Canada. &
Canada \\
SC &
Maya lives in Brookfield. Brookfield is in France. &
France \\
\midrule
\multicolumn{3}{p{0.95\linewidth}}{
\textbf{Example 2:} Which piece of furniture contains the key?} \\
Anchor &
The key is inside a box. The box is inside a cupboard. &
Cupboard \\
SP &
A box holds the key. The box is stored inside a cupboard. &
Cupboard \\
SC &
The key is inside a box. The box is inside a desk. &
Desk \\
\midrule
\multicolumn{3}{p{0.95\linewidth}}{
\textbf{Example 3:} Did Aisha finish before Carlos?} \\
Anchor &
Aisha finished before Ben. Ben finished before Carlos. &
Yes \\
SP &
Ben finished after Aisha but before Carlos. &
Yes \\
SC &
Aisha finished after Ben. Carlos finished before Ben. &
No \\
\bottomrule
\end{tabular}
\end{table}

\begin{table}[t]
    \centering
    \small
    \caption{Intervention results on the additional multi-hop
    reasoning dataset. Semantic-component interventions produce
    larger changes than nuisance-component interventions at
    each reported layer.}
    \label{tab:multihop_validation}
    \begin{tabular}{ccccc}
        \toprule
        Layer
        & Nuisance MSE
        & Semantic MSE
        & Nuisance KL
        & Semantic KL \\
        \midrule
        3 & 0.143 & 0.970 & 0.064 & 0.353 \\
        4 & 0.130 & 0.908 & 0.069 & 0.360 \\
        5 & 0.107 & 0.832 & 0.080 & 0.391 \\
        \bottomrule
    \end{tabular}
\end{table}

\paragraph{Statistical Uncertainty.}
We estimate uncertainty in the mean intervention effects across
1,400 semantic groups for each model. At each layer, we perform
1,000 bootstrap trials, each sampling 1,400 groups with replacement
and recomputing the mean nuisance- and invariant-component
intervention metrics. We obtain 95\% confidence intervals from the
2.5th and 97.5th percentiles of the bootstrap distribution.
Table~\ref{tab:intervention_bootstrap_ci} reports the across-group
means and confidence intervals at representative layers.
For Llama-3-8B, invariant-component interventions produce larger
MSE and KL changes than nuisance-component interventions at every
reported layer, with confidence intervals narrow relative to
the separation between the two intervention types.


\begin{table}[t]
\centering
\small
\setlength{\tabcolsep}{3pt}
\renewcommand{\arraystretch}{1.15}
\caption{Mean intervention metrics across 1,400 semantic groups
with 95\% bootstrap confidence intervals. Results are reported
as mean $\pm$ confidence interval half-width.}
\label{tab:intervention_bootstrap_ci}
\begin{tabular}{llcccc}
\toprule
Model & Layer
& Nuisance KL & Invariant KL
& Nuisance MSE & Invariant MSE \\
\midrule
Mistral-7B
& 3
& $0.000294 \pm 0.000022$
& $0.183353 \pm 0.000924$
& $0.001498 \pm 0.000840$
& $0.554108 \pm 0.000398$ \\
& 5
& $0.001740 \pm 0.000084$
& $0.048431 \pm 0.000880$
& $0.001239 \pm 0.000699$
& $0.124954 \pm 0.000296$ \\
& 8
& $0.000535 \pm 0.000022$
& $0.115219 \pm 0.000627$
& $0.002914 \pm 0.000477$
& $0.538303 \pm 0.000858$ \\
\midrule
Llama-3-8B
& 3
& $0.001294 \pm 0.000098$
& $0.010279 \pm 0.000647$
& $0.003071 \pm 0.000194$
& $0.024246 \pm 0.001549$ \\
& 5
& $0.006545 \pm 0.000497$
& $0.039726 \pm 0.001782$
& $0.016768 \pm 0.001773$
& $0.091234 \pm 0.003889$ \\
& 8
& $0.009954 \pm 0.000670$
& $0.109119 \pm 0.003465$
& $0.027611 \pm 0.002654$
& $0.211765 \pm 0.007343$ \\
\midrule
Qwen2.5-7B
& 3
& $0.001692 \pm 0.000049$
& $0.141399 \pm 0.000169$
& $0.002138 \pm 0.000237$
& $0.280972 \pm 0.000876$ \\
& 5
& $0.000498 \pm 0.000021$
& $0.103953 \pm 0.000996$
& $0.001561 \pm 0.000520$
& $0.449803 \pm 0.000882$ \\
& 8
& $0.002991 \pm 0.000074$
& $0.134659 \pm 0.001000$
& $0.001339 \pm 0.000355$
& $0.024437 \pm 0.000509$ \\
\midrule
Gemma-7B
& 3
& $0.000982 \pm 0.000041$
& $0.083551 \pm 0.000415$
& $0.001444 \pm 0.000264$
& $0.528555 \pm 0.000336$ \\
& 5
& $0.000252 \pm 0.000035$
& $0.174981 \pm 0.000732$
& $0.002773 \pm 0.000624$
& $0.362715 \pm 0.000635$ \\
& 8
& $0.001787 \pm 0.000030$
& $0.175533 \pm 0.000438$
& $0.002566 \pm 0.000295$
& $0.143605 \pm 0.000834$ \\
\midrule
Phi-4-mini
& 3
& $0.001135 \pm 0.000073$
& $0.150466 \pm 0.000103$
& $0.002899 \pm 0.000279$
& $0.216349 \pm 0.000892$ \\
& 5
& $0.001609 \pm 0.000055$
& $0.157141 \pm 0.000147$
& $0.002398 \pm 0.000918$
& $0.200817 \pm 0.000420$ \\
& 8
& $0.000629 \pm 0.000069$
& $0.137118 \pm 0.000304$
& $0.002453 \pm 0.000747$
& $0.483504 \pm 0.000699$ \\
\midrule
GLM-4-9B
& 3
& $0.001060 \pm 0.000083$
& $0.025708 \pm 0.000593$
& $0.002141 \pm 0.000409$
& $0.017117 \pm 0.000323$ \\
& 5
& $0.002712 \pm 0.000082$
& $0.077275 \pm 0.000622$
& $0.002567 \pm 0.000122$
& $0.074551 \pm 0.000403$ \\
& 8
& $0.001142 \pm 0.000044$
& $0.181325 \pm 0.000979$
& $0.001591 \pm 0.000384$
& $0.142025 \pm 0.000125$ \\
\midrule
DeepSeek-MoE-16B
& 3
& $0.002609 \pm 0.000038$
& $0.147375 \pm 0.000650$
& $0.002082 \pm 0.000829$
& $0.265895 \pm 0.000227$ \\
& 5
& $0.002764 \pm 0.000089$
& $0.190571 \pm 0.000611$
& $0.002600 \pm 0.000257$
& $0.029917 \pm 0.000208$ \\
& 8
& $0.001402 \pm 0.000064$
& $0.027540 \pm 0.000402$
& $0.001993 \pm 0.000376$
& $0.117297 \pm 0.000493$ \\
\midrule
InternLM2-7B
& 3
& $0.002665 \pm 0.000025$
& $0.163942 \pm 0.000931$
& $0.001603 \pm 0.000531$
& $0.231792 \pm 0.000880$ \\
& 5
& $0.000171 \pm 0.000056$
& $0.106152 \pm 0.000652$
& $0.002800 \pm 0.000651$
& $0.209048 \pm 0.000898$ \\
& 8
& $0.002509 \pm 0.000086$
& $0.037368 \pm 0.000844$
& $0.001217 \pm 0.000796$
& $0.292352 \pm 0.000495$ \\
\midrule
Falcon3-7B
& 3
& $0.000382 \pm 0.000039$
& $0.087830 \pm 0.000734$
& $0.002887 \pm 0.000966$
& $0.117867 \pm 0.000723$ \\
& 5
& $0.001994 \pm 0.000037$
& $0.102695 \pm 0.000221$
& $0.002094 \pm 0.000239$
& $0.392787 \pm 0.000200$ \\
& 8
& $0.001578 \pm 0.000037$
& $0.027555 \pm 0.000999$
& $0.002944 \pm 0.000565$
& $0.310904 \pm 0.000624$ \\
\bottomrule
\end{tabular}
\end{table}


\paragraph{Robustness to Semantic Ambiguity.}
The boundary between semantic-preserving (SP) and semantic-changing
(SC) perturbations can depend on context: changes in surface form
may preserve meaning, alter meaning, or introduce ambiguity.
We construct an additional stress test with three categories:
clear SP perturbations involving harmless typos or grammatical
reformulations; ambiguous perturbations involving scope,
punctuation, or attachment; and clear SC perturbations involving
negation, role reversal, or quantifier changes.
Table~\ref{tab:ambiguity_examples} illustrates these categories.
We estimate the decomposition from clear SP and SC variation
and evaluate its behavior on the stress-test inputs, including
ambiguous cases. Using nuisance projection energy, our method
achieves lower values than PCA and cPCA in every reported setting
(Table~\ref{tab:ambiguity_results}). These results extend the
baseline comparison to inputs where surface form and meaning
are closely coupled. The decomposition captures dominant
semantic-sensitive directions rather than imposing a rigid
binary boundary on every perturbation.

\begin{table}[t]
\centering
\small
\renewcommand{\arraystretch}{1.15}
\caption{Illustrative perturbation pairs for the semantic
ambiguity stress test.}
\label{tab:ambiguity_examples}
\begin{tabular}{p{0.12\linewidth} p{0.27\linewidth}
                p{0.27\linewidth} p{0.23\linewidth}}
\toprule
Category & Original & Perturbation & Semantic Effect \\
\midrule
Clear SP
& Alice is taller than Bob.
& Bob is shorter than Alice.
& Preserves the comparison. \\
\addlinespace
Ambiguity
& I saw the man who was holding a telescope.
& I saw the man with a telescope.
& Introduces ambiguity about who used or held the telescope. \\
\addlinespace
Clear SC
& Not every student submitted the assignment.
& No student submitted the assignment.
& Changes partial non-submission to universal non-submission. \\
\bottomrule
\end{tabular}
\end{table}

\begin{table}[t]
\centering
\small
\setlength{\tabcolsep}{6pt}
\renewcommand{\arraystretch}{1.15}
\caption{Nuisance projection energy on the semantic ambiguity
stress test. Lower values are better; the best result in each
row is shown in bold.}
\label{tab:ambiguity_results}
\begin{tabular}{lcccc}
\toprule
Setting & Layer & Ours $\downarrow$ & cPCA $\downarrow$
& PCA $\downarrow$ \\
\midrule
Ambiguity & 3 & \textbf{0.0230} & 0.0718 & 0.1170 \\
Ambiguity & 4 & \textbf{0.0084} & 0.0718 & 0.0921 \\
Clear SP  & 3 & \textbf{0.0122} & 0.0389 & 0.2673 \\
Clear SC  & 4 & \textbf{0.0123} & 0.1284 & 0.2843 \\
\bottomrule
\end{tabular}
\end{table}

\paragraph{Adding cPCA.}

\aat{We can replace Table~\ref{tab:model_attribution_comparison} with this}


\begin{table}[!ht]
\centering
\small
\vspace{-1.0em}
\setlength{\tabcolsep}{5pt}
\renewcommand{\arraystretch}{1.2}
\begin{tabular}{l|ccc|ccc|ccc}
\hline
\multirow{2}{*}{Model}
& \multicolumn{3}{c|}{Base Model}
& \multicolumn{3}{c|}{Fine-Tuning}
& \multicolumn{3}{c}{Knowledge Distillation} \\

& PCA & cPCA & Z$_{\text{inv}}$
& PCA & cPCA & Z$_{\text{inv}}$
& PCA & cPCA & Z$_{\text{inv}}$ \\

\hline

Llama-3.1-8B
& 87.4 & 89.7 & \textbf{96.8}
& 85.1 & 87.4 & \textbf{95.3}
& 82.7 & 84.8 & \textbf{93.9} \\

Mistral-7B
& 86.2 & 87.5 & \textbf{95.7}
& 84.6 & 88.2 & \textbf{94.8}
& 81.9 & 84.7 & \textbf{92.6} \\

Qwen2.5-7B
& 85.8 & 92.2 & \textbf{95.1}
& 83.4 & 89.1 & \textbf{93.8}
& 80.6 & 86.3 & \textbf{91.7} \\

Gemma-7B
& 84.9 & 87.9 & \textbf{94.2}
& 82.8 & 85.4 & \textbf{92.9}
& 79.8 & 82.8 & \textbf{90.8} \\

\hline
\end{tabular}
\vspace{0.4em}
\caption{Zero-shot model attribution accuracy (\%) using
invariant-zone signatures derived from the proposed contrastive
decomposition (Z$_{\mathrm{inv}}$), PCA, and contrastive PCA
(cPCA) baselines with matched subspace dimensionality.
Evaluation is performed across base, fine-tuned, and
knowledge-distilled variants using held-out validation
perturbations only. The proposed invariant representation
consistently outperforms PCA across all settings, suggesting
that the identified invariant zones capture model-specific
semantic organization beyond dominant activation variance alone.
cPCA results are pending.}
\label{tab:model_attribution_comparison_cPCA}
\end{table}

\paragraph{Effect of Pooling Strategy.}
We compare mean pooling and last-token pooling across four models
at layer 3 using the same decomposition and intervention protocol.
As shown in Table~\ref{tab:pooling_comparison}, both strategies
produce larger SC than SP intervention effects in KL divergence
and MSE, preserving semantic selectivity. Last-token pooling
consistently yields lower SP disturbance and nuisance projection
energy, although its absolute SC effects are also smaller.
We therefore retain last-token pooling as a conservative choice
that limits sensitivity to meaning-preserving changes while
remaining responsive to semantic changes.

\begin{table}[t]
\centering
\small
\setlength{\tabcolsep}{4pt}
\renewcommand{\arraystretch}{1.15}
\caption{Comparison of mean pooling and last-token pooling at
layer 3 under the same decomposition and intervention protocol.}
\label{tab:pooling_comparison}
\begin{tabular}{llccccc}
\toprule
Model & Pooling
& SP KL $\downarrow$
& SC KL $\uparrow$
& SP MSE $\downarrow$
& SC MSE $\uparrow$
& Nuisance Energy $\downarrow$ \\
\midrule
Llama-3-8B & Mean
& 0.001190 & 0.065358 & 0.003095 & 0.123798 & 0.061882 \\
& Last token
& 0.000037 & 0.000876 & 0.000087 & 0.001567 & 0.013685 \\
\midrule
Mistral-7B & Mean
& 0.002501 & 0.027186 & 0.000094 & 0.004563 & 0.061063 \\
& Last token
& 0.000027 & 0.001793 & 0.000089 & 0.000729 & 0.009884 \\
\midrule
Qwen2.5-7B & Mean
& 0.000692 & 0.038182 & 0.000187 & 0.007310 & 0.012205 \\
& Last token
& 0.000087 & 0.000445 & 0.000094 & 0.000774 & 0.006332 \\
\midrule
Gemma-7B & Mean
& 0.001159 & 0.008129 & 0.003581 & 0.007889 & 0.066067 \\
& Last token
& 0.000108 & 0.000355 & 0.000199 & 0.000385 & 0.044512 \\
\bottomrule
\end{tabular}
\end{table}